\begin{abstract}
\noindent Sistemas de reconhecimento facial passaram, na última década, de tecnologia laboratorial a infraestrutura social, estando hoje presentes em desbloqueio de dispositivos, autenticação bancária, controle de fronteiras e identificação policial. Em escala industrial, a maior auditoria pública já realizada em biometria facial, o relatório NISTIR 8280 do National Institute of Standards and Technology, documentou diferenciais de dez a cem vezes na taxa de falso positivo entre grupos raciais em cento e oitenta e nove algoritmos comerciais. Apesar da consolidação técnica, o problema de \emph{fairness} persiste: o estado da arte atual em classificação racial multi-classe sobre o dataset FairFace, obtido com \emph{Vision-Language Models} \emph{fine-tuned}, atinge F1 macro de setenta e cinco por cento, porém com disparidade severa entre classes --- aproximadamente noventa por cento para faces Black e apenas sessenta por cento para faces Latinx/Hispanic, gap de trinta pontos percentuais estável através de múltiplas arquiteturas e datasets.

\noindent O balanceamento explícito de dados, primeira geração de resposta ao problema, mostra-se insuficiente para resolver essa disparidade. Trabalhos empíricos recentes indicam que parte substantiva do erro reflete heterogeneidade fenotípica intra-categorial que a rotulagem monolítica de raça no FairFace não captura, sustentada por evidência convergente de três disciplinas independentes --- antropologia biológica, genética populacional e sociologia identitária. A literatura sugere ainda que o tom de pele, e não a categoria racial nominal, é o \emph{driver} estrutural do gap.

\noindent Propõe-se, nesta dissertação, desenvolver e avaliar um pipeline de classificação racial em imagens faciais que incorpora o tom de pele explícito, na escala Monk Skin Tone de dez classes, como sinal auxiliar condicionante via mecanismo arquitetural \emph{Feature-wise Linear Modulation} sobre \emph{backbone} ConvNeXt-T. O pipeline é organizado em seis etapas encadeadas: auditoria fenotípica do FairFace via um classificador MST pré-treinado, classificação racial condicionada por FiLM, comparação sistemática contra seis \emph{baselines} de mitigação --- FSCL+, Group DRO, FineFACE, \emph{adversarial debiasing} e dois controles arquiteturais ---, transferência \emph{fair} para tarefa \emph{downstream} de \emph{face recognition} sobre os datasets RFW e BFW, e síntese decompositiva final. A avaliação segue triangulação de métricas, com \emph{Disparity Ratio}, F1 da pior classe (\emph{worst-class} F1) e definições formais de \emph{Equal Opportunity} estratificadas por classe, arranjo desenhado para endereçar o Teorema da Impossibilidade em \emph{fairness} algorítmica e para permitir interpretação simultânea sob múltiplas definições formais de equidade.

\noindent Como contribuições esperadas, delineiam-se três eixos: primeiro, a construção e disponibilização pública da primeira matriz de distribuição cruzada Monk Skin Tone versus classes raciais sobre o \emph{validation set} do FairFace; segundo, a primeira instância empírica documentada do uso de tom de pele como contexto arquitetural em \emph{fairness} facial \emph{multi-class}; e terceiro, uma decomposição quantitativa do erro de classificação Latinx em componente fenotípico irredutível --- derivado do \emph{overlap} MST intra-categoria --- e componente algorítmico mitigável, oferecendo diagnóstico estrutural no lugar de simples redução agregada de F1 macro. O trabalho é conduzido em conformidade com o \emph{European AI Act} de 2024, que estabelece obrigatoriedade de auditoria de \emph{fairness} em sistemas biométricos classificados como de alto risco, e observa integralmente as diretrizes éticas institucionais aplicáveis à pesquisa computacional sobre datasets secundários publicamente disponíveis.

\medskip
\noindent\textbf{Palavras-chave:} reconhecimento facial; equidade algorítmica; viés racial; Monk Skin Tone; Feature-wise Linear Modulation; FairFace.
\end{abstract}

\vspace{1em}

\begin{abstract}
\noindent[English abstract to be added before submission --- mirror the Portuguese version.]
\end{abstract}

\newpage
